\section{Trainer Specification}
\label{supp:trainer}

\textit{Purpose of this supplement.}\;
This supplement records the full hyperparameter configuration,
learning-rate schedule, negative-sampling policy, and RNG seeding
protocol shared by all 310 main-factorial training runs.

\subsection{Hyperparameters}

All 310 main-factorial training runs share a common configuration.
Optimisation uses AdamW with linear decay over each stage's update
budget, learning rate $2\times 10^{-5}$, gradient clipping at L2
norm 1.0, and FP32 precision throughout. Each optimiser step
processes one mini-batch of 4 examples: 1 positive sampled from the
active shard pool and 3 negatives appended by online sampling (a 3:1
negative-to-positive ratio within each step). Each training stage
runs for 2{,}048 optimiser updates with evaluation every 64 updates
on a seeded 12\% in-training development split; full-vocabulary
macro-F1 drives early stopping with patience 10 evaluations and a
256-step minimum.

\begin{table}[H]
\centering
\caption{Full hyperparameter inventory.}
\label{tab:s4-hyperparams}
\small
\begin{tabular}{@{}lll@{}}
\toprule
\textbf{Group} & \textbf{Parameter} & \textbf{Value} \\
\midrule
Schedule    & First-stage updates                         & 2{,}048 \\
            & Second-stage updates (staged only)          & 2{,}048 \\
            & Optimiser/scheduler reset between stages    & yes \\
\addlinespace
Optimiser   & Optimiser                                   & AdamW \\
            & Learning rate                               & $2\times 10^{-5}$ \\
            & LR schedule                                 & linear decay (no warmup) \\
            & Gradient clipping (L2)                      & 1.0 \\
            & Mixed precision                             & disabled (FP32) \\
\addlinespace
Evaluation  & \code{eval\_every\_steps}                   & 64 \\
            & Selection metric                            & dev macro-F1 (full vocab) \\
            & Early-stopping patience                     & 10 evaluations \\
            & Minimum updates before stopping             & 256 \\
\addlinespace
Sampling    & Mini-batch size                             & 4 total (1 positive + 3 negatives) \\
            & Negative-to-positive ratio (per step)       & 3:1 \\
            & Maximum candidate negatives per positive    & 64 \\
\bottomrule
\end{tabular}
\end{table}

\subsection{Budget-pass arithmetic}

Each stage's 2{,}048 optimiser updates consume 2{,}048 positives
(one positive per step). A \emph{budget pass} is one full traversal
of the capped pool, so the number of passes equals positives consumed
divided by pool size. The single-corpus schedule (T1B) trains on
1{,}760 BioRED positives ($2{,}048/1{,}760\approx 1.16$ budget
passes); the multi-corpus flat schedule (T1F) draws from 5{,}280
pooled positives ($2{,}048/5{,}280\approx 0.39$ budget passes); and
the staged second stage draws from 4{,}456 oncology-projected
positives. The training regime is a deliberately fixed-budget
exposure to a shard-capped multi-source pool, not a full epoch over
all available relations; this isolates schedule effects from
training-set-size effects across encoders.

\subsection{Random-number control}

Each training run is keyed by a single integer seed in $\{1, \ldots, 20\}$.
The seed is propagated to Python's \code{random}, NumPy's RNG, the
CPU PyTorch RNG, the CUDA RNG, the development-split sampler, and
the online negative sampler (offset by $+101$ for the first stage and
$+202$ for the second). RNG-state snapshots are saved at every
checkpoint, so the exact batch stream from any checkpoint forward
can be reproduced.